\chapter{知识的共享}
\label{chap:ket-sharing}

花卉识别系统现在需要同时回答三个问题：照片中的花属于什么类别，是否出现病害，以及目标位于哪里。分别训练三个模型，容易重复学习边缘、纹理与形状；全部共用一个网络，又可能让任务差异被过早抹去。类别识别可以忽略的局部斑点，恰好可能是病害判断的关键。共同使用数据和参数只是起点，真正的问题是：哪些经验值得共同使用，哪些差异必须保留，以及某个任务的改善是否以其他任务受损为代价。

\term{多任务学习}（multi-task learning, MTL）让多个任务的训练过程相互联系，使一个任务获得的经验能够影响其他任务的学习。本章固定三个常驻任务$\mathcal T_1,\mathcal T_2,\mathcal T_3$，依次表示花卉类别、病害判断和对象定位。任务$k$具有分布$\mathcal P_k$、预测器$f_k$和损失$\ell_k$，其风险为
\begin{equation}
  R_k(f_k)=
  \mathbb E_{(\bm{x},y_k)\sim\mathcal P_k}
  [\ell_k(f_k(\bm{x}),y_k)].
  \label{eq:ket-sharing-task-risk}
\end{equation}
这里$y_k$可以是类别、病害状态或位置标注，不要求三个任务具有相同的输出空间。训练集记为$S_{\mathrm{train},k}$，经验风险为$\hat R_k$；后文用$L_k$表示当前批次上的任务损失。训练损失的变化提供局部优化证据，独立评价集上的风险与指标才用于判断共享收益。

本章依次考察任务组织、共享位置、联合目标和训练资源四项设计。它们并非四种互斥算法：一个系统可以让部分任务共享专家，用任务关系正则联系私有参数，再以动态权重协调训练。讨论每项设计时，应暂时固定其他条件，才能辨认改变来自哪里。第~\ref{chap:ket-transfer-adaptation}章研究已有经验怎样服务于新目标；这里研究多个当前任务怎样共同形成经验。把分别取得的模型再组合或融合，则属于第~\ref{chap:ket-integration}章的知识整合问题。

共享设计至少有两组正交坐标。结构侧用“共享载体 × 耦合阶段”描述：载体可以是输入样本、共同表示与参数、参数先验或输出关系，耦合可以发生在数据组织、前向计算、训练约束或部署调用。优化侧用“偏好目标 × 随机估计 × 梯度算子”描述：任务权重规定想优化什么，采样和缺失标签处理决定怎样估计它，PCGrad、MGDA等算子再决定如何把已取得的任务梯度变成更新。改变后一项不能自动修复前两项的偏差。

\section{共享任务的组织}
\label{sec:ket-sharing-organization}

三个任务都使用花卉照片，并不意味着它们必须全部共用一个模型。组织方案需要回答三个不同的问题：哪些预测要求值得纳入训练，任务之间的影响怎样估计，以及依据这些影响允许哪些任务共同学习。主任务与辅助任务描述收益要求，全体、任务组与个体描述共享层级，不能将两者当作同一条分类轴。组织方案也可以修订；试训发现病害任务持续受损时，应回头检查分组，而非把它当作不可改变的前提。

\subsection{候选任务与收益要求的确定}
\label{sec:ket-sharing-candidates}

从部署要求出发，类别、病害和定位若都要交付，就都属于需要独立验收的任务。三者不必具有相同权重：病害漏检的后果可能更严重，定位也可能只需达到后续裁剪可用的精度。应先规定每项任务的最低要求、允许的退化幅度和总体目标，再讨论怎样设置损失权重。以单任务基线风险$R_k^{\mathrm{single}}$为参照，沿用式~\eqref{eq:ket-reuse-risk-gain}中正值表示改善的方向，一种收益约定是
\begin{equation}
  \Delta_k=R_k^{\mathrm{single}}-R_k^{\mathrm{joint}}
  \geq-\epsilon_k,
  \qquad k\in C,
  \label{eq:ket-sharing-benefit-constraint}
\end{equation}
其中$C$是必须交付的任务集合，$\epsilon_k\geq0$是事先约定的容许退化量。$\Delta_k>0$表示共享降低了风险，$\Delta_k<0$表示该任务受损；$\epsilon_k=0$时不允许风险增加。不同任务的风险单位可以不同，因此不宜直接把这些收益相加，跨指标比较见第~\ref{sec:ket-evaluation-sharing}节。实际验收还要考虑估计误差；该式表达设计要求，不是一次训练即可保证的结论。

如果只需部署类别识别，可以增加花瓣颜色、叶缘形状等属性预测作为\term{辅助任务}（auxiliary task），即只在训练中提供额外要求、最终以主任务收益判断价值的任务。属性标签为中间表示保留某些可解释局部信息提供训练压力，主任务因而可能不必仅凭有限类别标签发现这些信息。Caruana的多任务学习工作用共同隐藏表示说明了这种间接影响：各任务可以使用不同的输出单元，却共同改变一部分内部表示\scite{caruana1997}

这种影响并不总是有益。花色有助于识别某些物种，但若主任务需要跨季节稳定识别，而辅助任务要求精细区分光照造成的色差，它也可能占用有限容量。辅助任务自身损失下降，不能替代主任务验证。合理的退出条件包括主任务收益在多个训练重复中不稳定、辅助标签费用超过收益，或加入辅助任务后无法满足推理预算。只交付主任务时，可以在部署中删除不再需要的辅助头；如果主任务前向计算依赖辅助预测，则不能直接删除这条路径。

额外预测要求也可以由已有数据构造，而不必每项都重新人工标注。Ando和Zhang研究了从多个预测问题中学习共同结构：多个任务共同限制可选的表示或假设空间，而任务特有函数仍分别拟合自己的目标\scite{ket-ando2005-structure}辅助目标怎样生成、核验以及是否泄漏答案，属于第~\ref{part:efficient-knowledge-use}部分的知识获取主题。本节只确定它是否值得成为候选，以及采用它必须满足什么收益条件；是否真正共享还需训练中的关系证据。

\subsection{任务间共享关系的估计}
\label{sec:ket-sharing-affinity}

语义关系给出尝试共享的理由，不能直接给出共享收益。类别与病害标签可能相关，但相关性也可能来自某个植物园恰好只拍摄了患病的某些物种。类别任务若利用这种捷径，可能损害病害模型在其他物种上的判断。需要测量的是：在指定结构与训练状态下，一个任务引起的更新怎样改变另一个任务的损失。

设$\bm{\theta}_{\mathrm{s}}^{(t)}$为阶段$t$的共享参数，$\bm{\theta}_k^{(t)}$为任务私有参数。固定当前批次及全部私有参数，对任务$k$做一次只涉及共享块的虚拟更新：
\begin{align}
  \bm g_k^{(t)}
    &=\nabla_{\bm{\theta}_{\mathrm{s}}}L_k^{(t)},
  \\
  \widetilde{\bm{\theta}}_{\mathrm{s},k}^{(t)}
    &=\bm{\theta}_{\mathrm{s}}^{(t)}-\eta\bm g_k^{(t)}.
  \label{eq:ket-sharing-virtual-step}
\end{align}
“虚拟”指测量后恢复原状态，不把依次测量多个任务变成连续真实训练。令$L_{j\mid k}^{(t)}$表示共享块换成$\widetilde{\bm{\theta}}_{\mathrm{s},k}^{(t)}$后任务$j$的损失。在$L_j^{(t)}>0$时，用有向的\term{任务亲和性}（task affinity）衡量这次更新对另一个任务的相对影响：
\begin{equation}
  A_{k\to j}^{(t)}
  =1-\frac{L_{j\mid k}^{(t)}}{L_j^{(t)}}.
  \label{eq:ket-sharing-affinity}
\end{equation}
正值表示任务$k$的这次更新降低了任务$j$的损失，负值表示升高。Fifty等据此汇总训练过程中的任务亲和性，用来寻找共同训练的任务组\scite{ket-fifty2021-affinity}

实际估计时，每个方向都应使用可比较的批次、相同的步长和共享层范围，跨多个批次与阶段记录均值及波动。当前损失接近零时，比例会很不稳定，可同时报告绝对变化或另设明确的归一化尺度。用训练批次估计得到的是训练中的局部影响；若用验证数据挑选分组，这些数据就承担模型选择职责，不能再作为最终测试集。

在共享块梯度局部Lipschitz连续的条件下，一阶展开解释了这种测量与梯度的关系：
\begin{equation}
  L_{j\mid k}^{(t)}-L_j^{(t)}
  =-\eta(\bm g_j^{(t)})^{\top}\bm g_k^{(t)}
  +O(\eta^2).
  \label{eq:ket-sharing-local-effect}
\end{equation}
小步长下，正内积预示相互帮助，负内积预示局部冲突。不过，原始梯度内积是对称的，有限步长下的损失变化却受到任务曲率影响；式~\eqref{eq:ket-sharing-affinity}还使用不同的损失作分母。不能由内积对称推出实际亲和性对称。

下面的数值仅为教学构造。用一个共享标量$\theta$代表某种特征强度，令类别损失为$L_{\text{类}}(\theta)=\tfrac12(\theta-1)^2$，属性损失为$L_{\text{属}}(\theta)=5(\theta-0.1)^2$。在$\theta=0$处，两者梯度都是$-1$。取$\eta=0.3$，无论由哪项任务提议，虚拟参数都变成$0.3$；类别损失从$0.5$降到$0.245$，属性损失却从$0.05$升到$0.2$。因此$A_{\text{属}\to\text{类}}=0.51$，而$A_{\text{类}\to\text{属}}=-3$。单向帮助来自属性损失较大的曲率与跨过最优点的有限步长，不是来自不对称的梯度内积。改用足够小的步长，例如$0.01$，两项损失都会下降。

同样在教学构造中，若两项损失均为正曲率、最优点都在当前位置右侧，足够小的更新可以相互帮助；若病害损失为$\tfrac12(\theta+1)^2$，其在零点的梯度为$1$，就与上述类别梯度局部冲突。这些例子说明的是相互影响的可能形态，不能据此断言真实花卉任务永久属于某种关系。表示、标签噪声、训练阶段乃至共享层数改变，都可能改变估计。

第~\ref{sec:ket-reuse-task-relations}节判断已有知识能否服务于目标。Taskonomy的迁移关系通过冻结来源表示、训练目标读出等操作考察这种可复用性；它不是两个任务同时改变共享参数时的关系\scite{ket-zamir2018-taskonomy}由类别表示容易读出属性，不等于属性损失加入训练后一定改善类别表示。对辅助任务，应重点看它对主任务的影响；对共同交付的三个任务，则必须分别保留帮助、受损和不确定的证据。

\subsection{任务分组与共享层级的确定}
\label{sec:ket-sharing-grouping}

有了关系证据，可以比较全部共同训练、全部分别训练以及分组训练。类别与定位若在当前骨干中较为相容，而病害需要保留另一组局部纹理，可先让前两者组成一组、病害单独训练，再与全部共享进行试验。分组并不只是在一张亲和性矩阵上画边：加入第三个任务会改变梯度合成、容量竞争和采样次数，两两都相容不保证三者共同训练仍然受益。

比较时必须固定预算。把一个共享网络拆成三个同样大的网络，总参数和推理计算通常都会增加；反过来，把三个小网络换成一个大骨干，也同时改变了容量。Standley等将任务分组与网络大小、数据规模及计算预算一起考察，说明共同训练关系不能脱离这些条件判断\scite{ket-standley2020-grouping}一个可执行的选择过程是：先用局部关系筛出少量候选分组，再按相同总计算预算配置各组模型，完整试训后依据各任务验证指标选择，最终只在独立测试集上报告一次结果。组合数量较多时，筛选能降低费用，但不能免去整组验证。

共享还可以发生在同一预测要求的多个地区副本之间。暂时只考虑花卉类别识别，用$a$索引植物园，$c(a)$表示植物园所属的地区组；此处的个体不再指类别、病害和定位三个预测要求。全体植物园可以共同学习花瓣边缘，地区组保留常见拍摄条件的差异，各植物园再保留局部校准。对同维且坐标对应的参数块，可写成
\begin{equation}
  \bm{\theta}_a
  =\bm{\mu}+\bm b_{c(a)}+\bm\delta_a,
  \label{eq:ket-sharing-hierarchy}
\end{equation}
其中$\bm{\mu}$表示全体共性，$\bm b_{c(a)}$和$\bm\delta_a$分别表示地区组与植物园的偏差。这个分解表达允许共享的层级，还没有规定怎样估计各项。地理相近的两个地区也可能使用完全不同的相机；因此，行政或地理层级可以提供候选结构，却不应直接当作最优任务关系。

数据不能汇集时，个性化需求更加明显。联邦多任务学习MOCHA保留不同节点的任务模型；沿用上面的个体索引，记为$\bm w_a$，并通过关系矩阵$\bm\Omega$联系它们。其基本形式是联合最小化各节点的经验损失与关系正则项，并交替处理任务参数和关系结构\scite{ket-smith2017-mocha}在关系固定时，节点利用本地数据近似求解局部子问题，再协调全局更新，而不是强迫所有节点具有一个完全相同的预测器。原方法的优化分析依赖凸任务损失、规定的正则条件以及对局部求解质量和节点掉线的限制，不能原样推广到任意非凸神经网络。数据留在本地也不等于模型更新不会泄漏信息；隐私要求回指第~\ref{chap:trust-privacy}章，通信协议和执行方式交给第~\ref{part:systems}部分。

异构任务与指令任务集合也需要这种组织工作。一种预测要求换用十个模板，不自然地产生十个独立任务；同一模板在不同语言中出现，也不意味着语言差异已经消失。应分别记录任务语义、数据来源、语言与模板，并确定评价以哪个单位汇总。组织阶段决定哪些任务之间允许交换经验，下面再决定经验在模型的什么位置发生交换。

\section{可共享知识的设计}
\label{sec:ket-sharing-structure}

即使三个任务已经分在同一组，仍有多种共享方式。它们可以共同执行一个特征映射，也可以保留独立模型、只用统计关系约束参数，还可以通过输出关系联系预测。前者共享计算，第二种共享参数的先验结构，第三种共享预测所需的关系。共同表示常由共同参数实现，因此不能把“共享表示”和“共享参数”直接当作两种平行且无关的方法。

共享之前先要对齐接口。类别、病害和定位可以观察同一图像，但必须确认图像裁剪、坐标原点、对象对应和标签含义一致。对语言任务，T5用文本到文本的统一接口，通过任务前缀与文本化答案表达不同预测要求；统一形式便于复用模型，但不会使分类与生成具有相同的难度或数据规模，其全部训练实验也不能一概解释为同时联合训练\scite{raffel2020t5}MultiModel则保留不同模态的输入输出模块，让图像、语音或文本通过各自接口进入共享内部计算\scite{ket-kaiser2017-multimodel}接口统一与接口分离都可以支持共享，关键是接口是否保留任务所需的信息，而不是是否把所有输出改写成同一种字符串。

\subsection{共同计算结构的设计}
\label{sec:ket-sharing-computation}

共同计算让多个损失改变同一映射。需要决定的是：哪些层始终被这些任务共同使用，哪些部件只在某个任务或某类输入出现时使用。共享路径与混合比例可以预先固定，也可以由任务身份、当前输入或二者共同选择。仅依赖任务身份的选择，在每个任务内部仍是固定的，因此固定与条件共享存在边界重叠；下面重点比较固定连接与逐输入混合，两者都可以保留私有分支。

\subsubsection{固定共享范围的设计}
\label{sec:ket-sharing-fixed}

最直接的结构是共同骨干加任务头：
\begin{equation}
  \bm z=h_{\bm{\theta}_{\mathrm{s}}}(\bm x),
  \qquad
  \hat y_k=q_{\bm{\theta}_k}(\bm z).
  \label{eq:ket-sharing-backbone}
\end{equation}
每个任务都通过自己的头解释$\bm z$，但所有任务对$\bm{\theta}_{\mathrm{s}}$的梯度共同塑造它。共享浅层意味着共同学习边缘等早期特征，较深的任务分支仍可选择不同空间细节；共享到深层则节省更多重复计算，却要求高层表示同时保留类别、病害与位置。类别识别可能偏好对位置变化不敏感的汇总表示，定位却需要空间对应，因此共享深度应围绕具体表示检查，而不是假定“越深越通用”。

线性模型能更清楚地显示共性与差异。令输入$\bm x\in\mathbb R^d$，共享基$\bm U\in\mathbb R^{d\times r}$，任务系数$\bm a_k\in\mathbb R^r$，则
\begin{equation}
  \hat y_k=\bm a_k^{\top}\bm U^{\top}\bm x,
  \qquad
  \bm W=\bm U\bm A,
  \label{eq:ket-sharing-subspace}
\end{equation}
其中$\bm A$的第$k$列是$\bm a_k$，$\bm W$的第$k$列是任务参数。所有任务在同一个低维子空间里选函数，但不必具有相同系数。多任务特征学习进一步通过共同基与系数的行稀疏性，鼓励任务使用一组较少的特征；这与直接预设秩$r$的教学形式相关，但并非完全相同的约束\scite{ket-argyriou2006-features}训练需要同时估计共同结构和任务系数，使用时只需执行选定表示与任务函数。维数过小会丢失任务差异，过大又可能失去通过共享限制估计复杂度的作用。

另一个选择是保留任务分支，在指定层交换激活。Cross-stitch网络对形状对齐的分支特征做线性混合。以两个任务在同一位置、同一通道上的激活为例，
\begin{equation}
  \begin{bmatrix}
    \widetilde h_1\\ \widetilde h_2
  \end{bmatrix}
  =
  \begin{bmatrix}
    a_{1,1}&a_{1,2}\\a_{2,1}&a_{2,2}
  \end{bmatrix}
  \begin{bmatrix}
    h_1\\h_2
  \end{bmatrix}.
  \label{eq:ket-sharing-cross-stitch}
\end{equation}
对角项控制保留本分支多少信息，非对角项控制引入另一分支多少信息。这些系数随其他参数一起训练，不必解释为概率，也不必强制每行和为1；上式只是两分支的局部计算，三个任务可以在选定分支间设置相应连接\scite{ket-misra2016-cross-stitch}

例如，教学系数两行为$(0.9,0.1)$和$(0.2,0.8)$，两个激活为$2,6$时，混合后分别为$2.4,5.2$。梯度既更新系数，也沿混合连接返回原分支。虽然系数是学出来的，它们在部署时不因当前图像重新计算；“允许在哪些层、哪些分支之间连接”的范围也是事先规定的。这与下一节的输入条件门控不同。

共享不足时，三个分支可能重复估计相同纹理；共享过多时，它们又可能被迫竞争有限特征。验证共享层数或混合位置时，需要控制总参数、特征维数与计算量。若某个分支比其他方案宽一倍，不能把它的改善全部归于共享方式。

\subsubsection{条件共享范围的设计}
\label{sec:ket-sharing-conditional}

同一个任务也未必始终需要同一组特征。病害判断处理清晰叶片时可能重视斑点，处理遮挡照片时则更依赖形状与上下文。条件共享让当前任务或输入参与选择，而不必为每种情况训练完全独立的模型。

\term{多门控混合专家模型}（Multi-gate Mixture-of-Experts, MMoE）使用$E$个共享专家$\bm e_m(\bm x)$，同时为每个任务设置门控网络。门控产生归一化权重，再交给任务头：
\begin{align}
  g_{k,m}(\bm x)
    &=\frac{\exp(a_{k,m}(\bm x))}
      {\sum_{r=1}^{E}\exp(a_{k,r}(\bm x))},
  \\
  \bm z_k(\bm x)
    &=\sum_{m=1}^{E}g_{k,m}(\bm x)\bm e_m(\bm x),
  \nonumber\\
  \hat y_k&=q_{\bm{\theta}_k}(\bm z_k(\bm x)).
  \label{eq:ket-sharing-mmoe}
\end{align}
专家提供可复用映射，任务特有门控决定本任务在这张图像上怎样混合，任务头再完成各自预测。MMoE据此让任务关系反映在不同门控对共同专家的使用中\scite{ket-ma2018-mmoe}

训练时，各任务损失同时更新本任务的头、门控以及它使用的专家；使用时，输入经过专家和所请求任务的门控，再得到任务输出。MMoE的混合同时依赖任务身份与输入：不同任务有不同门控，同一任务的权重也随图像变化。若所有任务共用一个只读输入的门控，同一图像就会得到相同的混合；Cross-stitch则沿用训练得到的固定系数。

共享专家之外，还可以保留任务私有分支或地区校准参数。它们为局部差异提供容量，却也增加需要估计的参数，不能单凭结构就断言个性化有效。门控及专家训练都受任务竞争影响：样本多且更新机会多的任务可能更早塑造专家。其他任务仍能更新专家，但可能难以在有限训练预算内充分纠正这种偏向。多个专家甚至可能学成近似相同的映射。

因此，检查条件共享时，应统计各任务及其输入子群上的平均门控权重、权重的集中程度，以及移除或替换某个专家后预测怎样变化，并与同等预算的固定共享对照。不同权重不自动意味着专家具有清楚的语义分工。标准MMoE采用稠密软混合，通常需要计算全部专家；只有进一步引入稀疏选择和相应执行机制，才可能减少每次激活的计算，具体预算在第~\ref{sec:ket-sharing-budget}节核算。这里的专家在共同训练中形成；如何组合已经取得的模块，回指第~\ref{sec:ket-integration-conditional-modules}节。

\subsection{共享先验与参数关系的设计}
\label{sec:ket-sharing-prior}

多个任务不共用一段前向计算，也可以共享统计信息。以下用$k$索引参与这种联系的任务，先要求$\bm w_k\in\mathbb R^d$同维且使用相同坐标：它可以是同一组特征上的线性预测参数，也可以是各任务网络中已经对齐、按相同顺序展开的同形参数块。对类别、病害和定位任务，可以只联系这样的编码器块，异构输出头各自保留，不直接相加或围绕同一个向量求平均；下式省略其余私有参数。

假设这些参数块围绕某个共同中心变化，数据少的任务便可以借助其他任务来估计合理范围，不必独自确定每个自由度。

一个简化的层次高斯模型为
\begin{equation}
  \bm\mu\sim\mathcal N(\bm0,\kappa^2\bm I),
  \qquad
  \bm w_k\mid\bm\mu\sim
  \mathcal N(\bm\mu,\tau^2\bm I).
  \label{eq:ket-sharing-hierarchical-prior}
\end{equation}
这里$\kappa,\tau>0$固定，$\tau$控制任务偏离共同中心的自由。若$\mathcal L_k(\bm w_k)$是任务数据的负对数似然，最大后验估计等价于最小化
\begin{equation}
  \sum_{k=1}^{K}\mathcal L_k(\bm w_k)
  +\frac{1}{2\tau^2}\sum_{k=1}^{K}
     \|\bm w_k-\bm\mu\|_2^2
  +\frac{1}{2\kappa^2}\|\bm\mu\|_2^2.
  \label{eq:ket-sharing-hierarchical-map}
\end{equation}
因此，先验成为把任务拉向共同中心的正则项。给定各$\bm w_k$，中心的最优解为
\begin{equation}
  \bm\mu
  =\frac{\sum_{k=1}^{K}\bm w_k}
    {K+\tau^2/\kappa^2}.
\end{equation}
给定中心时，各任务再在拟合自身数据与偏离中心之间取舍。数据量小、似然提供约束弱的任务会受到更强的相对收缩；这不等于把每个任务的参数直接设为平均值。

如果参与联系的任务具有多个群组，可以为任务引入簇身份$c_k$，令$\bm w_k\mid c_k$围绕相应的$\bm\mu_{c_k}$变化，再估计任务属于各簇的概率。Bakker和Heskes以贝叶斯多任务模型研究共同先验、任务聚类与门控，使相近任务能够通过簇结构共享信息\scite{ket-bakker2003-clustering}式~\eqref{eq:ket-sharing-hierarchical-prior}是解释收缩作用的简化模型，不是该论文全部推断过程。实际可以事先按可靠知识指定簇，也可以从数据估计；后者增加了灵活性，同时需要足够证据辨认关系。

共同中心主要表达“接近”。\term{任务关系学习}（Multi-Task Relationship Learning, MTRL）进一步让不同参数方向上的任务变化具有协方差结构。令$\bm W=[\bm w_1,\ldots,\bm w_K]\in\mathbb R^{d\times K}$，这里采用的矩阵正态先验约定各行相互独立，且每行都是零均值、协方差为$\bm\Omega\succ0$的$K$维高斯向量。其负对数密度包含
\begin{equation}
  \frac12\operatorname{tr}
    (\bm W\bm\Omega^{-1}\bm W^{\top})
  +\frac d2\log|\bm\Omega|.
  \label{eq:ket-sharing-matrix-prior}
\end{equation}
当$\bm\Omega$也要估计时，行列式项不能随意丢弃。MTRL为得到相应的凸正则化形式，以迹约束控制关系矩阵的尺度，构造
\begin{align}
  \min_{\bm W,\bm\Omega}\quad&
  \sum_{k=1}^{K}\hat R_k(\bm w_k)
  +\frac{\lambda_1}{2}\|\bm W\|_{\mathrm F}^2
  \nonumber\\
  &+\frac{\lambda_2}{2}
    \operatorname{tr}(\bm W\bm\Omega^{-1}\bm W^{\top})
  \nonumber\\
  \text{满足}\quad&
  \bm\Omega\succeq0,
  \qquad \operatorname{tr}(\bm\Omega)=1.
  \label{eq:ket-sharing-mtrl}
\end{align}
这一步改变了估计准则，并不是原高斯最大后验目标的代数恒等改写\scite{ket-zhang2010-mtrl}

关系固定且损失为凸函数时，任务参数子问题是带耦合正则的凸优化；参数固定后，可更新任务关系。令$\bm B=\bm W^{\top}\bm W$，若$\bm W\ne\bm0$，对应的半正定最优解为
\begin{equation}
  \bm\Omega^\star
  =\frac{\bm B^{1/2}}
    {\operatorname{tr}(\bm B^{1/2})}.
  \label{eq:ket-sharing-mtrl-omega}
\end{equation}
可以从$\bm B$的特征分解理解它：若特征值为$b_j$，对齐特征方向后需在$\sum_j\omega_j=1$下最小化$\sum_j b_j/\omega_j$，拉格朗日条件给出$\omega_j\propto\sqrt{b_j}$。交替拟合$\bm W$和更新$\bm\Omega$，就让当前多个任务共同确定参数关系；部署时仍可各自使用$\bm w_k$预测，不必保留其他任务的前向网络。

式~\eqref{eq:ket-sharing-mtrl}中的逆矩阵写法在$\bm\Omega\succ0$时直接成立。在半正定边界上，应使用相应的伪逆，并要求$\operatorname{range}(\bm W^{\top})$包含于$\operatorname{range}(\bm\Omega)$，否则将正则值视为无穷大；数值实现也可加入正定下界，但此时闭式更新需相应调整。若$\bm W=\bm0$，迹为1的关系矩阵都得到同样的零关系正则，式~\eqref{eq:ket-sharing-mtrl-omega}不能作零除。这些条件防止把低秩解误当成普通可逆矩阵。

负相关也能承载有用关系。对两个使用相同特征坐标的任务，取
\begin{equation}
  \bm\Omega=\frac12
  \begin{bmatrix}1&\rho\\\rho&1\end{bmatrix},
  \qquad |\rho|<1.
\end{equation}
它对应的关系正则为
\begin{equation}
  \operatorname{tr}(\bm W\bm\Omega^{-1}\bm W^{\top})
  =
  \frac{2\bigl(\|\bm w_1\|_2^2+\|\bm w_2\|_2^2
  -2\rho\bm w_1^{\top}\bm w_2\bigr)}
  {1-\rho^2}.
  \label{eq:ket-sharing-negative-correlation}
\end{equation}
若$\bm w_2\approx-\bm w_1$，负的$\rho$反而给予这种相反趋势较小的惩罚。两个任务可以分别预测某种观测条件下互补的状态，一项参数的估计便能帮助预测另一项的相反变化。这里的负相关是参数统计关系，负迁移则是共同学习使任务表现变差，二者不是同一概念。

这种解释依赖坐标与参数化。如果各任务的特征基没有对齐，单独翻转一个任务的某个特征符号，或重排其隐藏单元，就可能改变表观参数关系而不改变预测函数；不能把这种参数差异直接当作任务差异。协方差既不是固定的语义标签，也不是因果关系。共同先验形状错误、簇划分错误或关系估计样本不足，都可能把任务收缩到不合适的位置。第~\ref{sec:ket-transfer-prior}节在新任务到来前准备可迁移先验；这里则让多个当前任务共同估计并使用它，收益应在这些任务上分别检验。

\subsection{共享输出关系与约束的设计}
\label{sec:ket-sharing-outputs}

参数关系不易解释时，输出之间可能具有更直接的联系。例如，在确认是同一对象、裁剪坐标一致的条件下，定位结果可以告诉病害模型应关注哪个区域；类别与可见属性也可能具有相容关系。关系必须带条件：照片中有多朵花时，左侧对象的类别不能约束右侧对象的病害；看不到叶片时，“未预测出叶片斑点”也不能直接推出健康。

一种做法是把辅助预测变成后续目标的输入。PAD-Net先从图像产生多个中间任务的预测，再把这些预测转换成特征，通过拼接、消息传递或注意力等形式支持最终目标。原论文的中间任务包括深度、表面法向、轮廓和语义，最终目标是深度与语义预测；训练和测试都以图像为输入\scite{ket-xu2018-padnet}

对应到花卉系统，可以从共同编码器$h$和任务映射$t_k$构造如下教学结构。这是一个独立的消息传递例子，不默认接入前面的MMoE，也不是PAD-Net的原实验：
\begin{align}
  \bm z&=h(\bm x),
  \qquad \bm z_k=t_k(\bm z),
  \nonumber\\
  \widetilde y_j&=u_j(\bm z),
  \qquad
  \bm r_j=\phi_j(\widetilde y_j),
  \nonumber\\
  \bm v_k&=\bm z_k+
     \sum_{j\ne k}
     \bm a_{k,j}(\bm r_k,\bm r_j)
       \odot\psi_{k,j}(\bm r_j),
  \nonumber\\
  \hat y_k&=q_k(\bm v_k).
  \label{eq:ket-sharing-output-message}
\end{align}
其中$t_k$形成供任务$k$接收消息的基础表示，$\phi_j$把中间预测转为特征。若接收分支的高、宽和通道数为$H_k,W_k,C_k$，则$\bm z_k,\bm v_k\in\mathbb R^{H_k\times W_k\times C_k}$；$\psi_{k,j}(\bm r_j)$和门控输出$\bm a_{k,j}(\bm r_k,\bm r_j)$也须具有这一形状，才能逐元素相乘并与$\bm z_k$相加。$\psi_{k,j}$负责通道映射与空间对齐；图像级类别概率可映射后广播到接收网格，定位热图则需按相同裁剪坐标重采样。

门控调节各位置接收多少信息，$\odot$表示逐元素相乘。中间定位热图可以帮助病害分支聚焦对象区域，类别概率也可以提供物种相关线索。训练同时约束中间预测和最终输出，梯度沿这些可微路径返回；部署时由同一图像重新生成中间预测，再执行后续预测。

真实辅助标签只负责训练监督，不能悄悄替代部署时的$\widetilde y_j$。若训练用真实区域、测试用预测区域，后续任务面对的输入误差条件就发生了变化。传递完整概率图通常比过早取一个硬标签保留更多不确定性，但仍可能传播错误；应测试区域偏移、错误类别和低置信输入下的后续表现。注意力权重也不自动等于校准后的置信概率。

另一种做法不增加输出之间的前向依赖，只在训练目标中约束一致性。设病害分支给出同一对象的像素病变概率$p_{\text{病}}(u\mid\bm x)$，定位分支给出对象区域概率$p_{\text{区}}(u\mid\bm x)$。若标签定义明确规定所预测的病变应位于该对象区域内，可加入教学约束
\begin{equation}
  L_{\mathrm{rel}}
  =\frac{1}{|\mathcal U|}
    \sum_{u\in\mathcal U}
    p_{\text{病}}(u\mid\bm x)
    \bigl[1-p_{\text{区}}(u\mid\bm x)\bigr],
  \label{eq:ket-sharing-output-constraint}
\end{equation}
其中$\mathcal U$为对齐的像素集合。联合损失为原有监督项加$\lambda_{\mathrm{rel}}L_{\mathrm{rel}}$。这里为解释空间约束，病害任务额外具有病变位置输出；若只输出图像级有无病害，不能直接使用该式。没有监督锚定时，病变概率全部变成零或区域全部变成一都能减小此项，因此它不能替代任务标签。约束错误或权重过大也会伤害本来正确的预测。

预测传递把关系放进前向计算，需要承担错误传播与推理依赖；一致性正则把关系放进训练目标，部署时可以保留独立输出，但无法自动保证每次预测严格满足约束。若各模块已训练完成，只需按接口串接计算，则是第~\ref{sec:ket-integration-fixed-modules}节的组合问题；这里还考察这些关系如何约束共同学习。规则怎样获得与核验、多来源答案怎样形成训练依据，属于知识获取主题。图~\ref{fig:ket-sharing-locations}概括了三种共享位置，并展开条件门控。它们可以组合使用，但图中的实线和虚线表达不同关系，不能都读成一条前向计算链。

\begin{figure*}[htbp]
  \centering
  % Three sharing locations and a conditional-computation inset.
% Conceptual diagram only; all dimensions are in millimetres.
\begin{tikzpicture}[
  x=1mm,y=1mm,
  font=\figurefont\footnotesize,
  text=BookInk,
  >={Latex[length=1.8mm,width=1.3mm]},
  panel/.style={draw=BookRule,line width=0.6pt},
  model/.style={draw=BookTeal,fill=FigureModel!10,line width=0.8pt,
    minimum height=10mm,text width=11mm,align=center,inner sep=1mm},
  prior/.style={draw=BookBlue,fill=FigureInput!10,line width=0.8pt,
    minimum height=11mm,text width=27mm,align=center,inner sep=1mm},
  output/.style={text width=14mm,align=center,inner sep=0pt},
  relation/.style={draw=BookAmber,fill=FigureLoss!10,line width=0.8pt,
    minimum height=10mm,text width=40mm,align=center,inner sep=1mm},
  flow/.style={draw=BookInk,line width=0.8pt,-{Latex}},
  wire/.style={draw=BookInk,line width=0.8pt},
  stat/.style={draw=BookBlue,dashed,line width=0.8pt},
  constraint/.style={draw=BookAmber,dashed,line width=0.8pt},
  heading/.style={anchor=west,font=\figurefont\small}
]
  \path[use as bounding box] (0,5) rectangle (169,-160);
  \node[heading] at (1,0) {(a) 表示/参数 × 前向计算};
  \node[heading] at (59,0) {(b) 参数 × 训练先验};
  \node[heading] at (117,0) {(c) 输出 × 训练约束};
  \draw[panel] (0,-7) rectangle (53,-74);
  \draw[panel] (58,-7) rectangle (111,-74);
  \draw[panel] (116,-7) rectangle (169,-74);

  \node[model,text width=27mm] (shared) at (26.5,-20)
    {共同映射\\$h_{\bm{\theta}_{\mathrm s}}$};
  \node[model] (head1) at (8.5,-43) {$q_1$};
  \node[model] (head2) at (26.5,-43) {$q_2$};
  \node[model] (head3) at (44.5,-43) {$q_3$};
  \node[output] (outa1) at (8.5,-62) {类别\\$\hat y_1$};
  \node[output] (outa2) at (26.5,-62) {病害\\$\hat y_2$};
  \node[output] (outa3) at (44.5,-62) {定位\\$\hat y_3$};
  \draw[wire] (shared.south) -- (26.5,-31);
  \draw[flow] (26.5,-31) -| (head1.north);
  \draw[flow] (26.5,-31) -- (head2.north);
  \draw[flow] (26.5,-31) -| (head3.north);
  \draw[flow] (head1.south) -- (outa1.north);
  \draw[flow] (head2.south) -- (outa2.north);
  \draw[flow] (head3.south) -- (outa3.north);

  \node[prior] (prior) at (84.5,-20)
    {共同先验\\$\bm\mu,\bm\Omega$};
  \node[model] (modelb1) at (66.5,-43) {$f_{\bm w_1}$};
  \node[model] (modelb2) at (84.5,-43) {$f_{\bm w_2}$};
  \node[model] (modelb3) at (102.5,-43) {$f_{\bm w_3}$};
  \node[output] (outb1) at (66.5,-62) {类别\\$\hat y_1$};
  \node[output] (outb2) at (84.5,-62) {病害\\$\hat y_2$};
  \node[output] (outb3) at (102.5,-62) {定位\\$\hat y_3$};
  \draw[stat] (prior.south) -- (84.5,-31);
  \draw[stat] (84.5,-31) -| (modelb1.north);
  \draw[stat] (84.5,-31) -- (modelb2.north);
  \draw[stat] (84.5,-31) -| (modelb3.north);
  \draw[flow] (modelb1.south) -- (outb1.north);
  \draw[flow] (modelb2.south) -- (outb2.north);
  \draw[flow] (modelb3.south) -- (outb3.north);

  \node[model] (modelc1) at (124.5,-20) {$f_1$};
  \node[model] (modelc2) at (142.5,-20) {$f_2$};
  \node[model] (modelc3) at (160.5,-20) {$f_3$};
  \node[output] (outc1) at (124.5,-43) {类别\\$\hat y_1$};
  \node[output] (outc2) at (142.5,-43) {病害\\$\hat y_2$};
  \node[output] (outc3) at (160.5,-43) {定位\\$\hat y_3$};
  \node[relation] (relation) at (142.5,-63)
    {同对象输出约束\\$L_{\mathrm{rel}}$};
  \draw[flow] (modelc1.south) -- (outc1.north);
  \draw[flow] (modelc2.south) -- (outc2.north);
  \draw[flow] (modelc3.south) -- (outc3.north);
  \draw[constraint] (outc1.south) -- (124.5,-53) -- ($(relation.north)+(-10,0)$);
  \draw[constraint] (outc2.south) -- (relation.north);
  \draw[constraint] (outc3.south) -- (160.5,-53) -- ($(relation.north)+(10,0)$);

  \node[heading] at (0,-85) {(d) 表示/参数 × 条件前向：专家由各任务共同训练};
  \node[prior,text width=16mm] (input) at (12,-116)
    {同一图像\\$\bm x$};
  \node[model,text width=25mm] (expert1) at (54,-101)
    {专家1：$\bm e_1(\bm x)$};
  \node[model,text width=25mm] (expert2) at (54,-116)
    {专家2：$\bm e_2(\bm x)$};
  \node[model,text width=25mm] (expert3) at (54,-131)
    {专家3：$\bm e_3(\bm x)$};
  \node[prior,text width=28mm] (gate) at (54,-151)
    {任务$k$的门控\\$\bm g_k(\bm x)$};
  \node[model,text width=24mm,minimum height=16mm] (mix) at (111,-116)
    {按权重混合\\$\bm z_k$};
  \node[model,text width=24mm,minimum height=16mm] (task) at (151,-116)
    {任务头$q_k$\\输出$\hat y_k$};
  \node[output,text width=30mm] at (151,-134)
    {$k=1,2,3$\\类别、病害、定位};
  \draw[wire] (input.east) -- (28,-116);
  \draw[flow] (28,-116) |- (expert1.west);
  \draw[flow] (28,-116) -- (expert2.west);
  \draw[flow] (28,-116) |- (expert3.west);
  \draw[flow] (28,-116) |- (gate.west);
  \draw[flow] (expert1.east) -- (mix.north west);
  \draw[flow] (expert2.east) -- (mix.west);
  \draw[flow] (expert3.east) -- (mix.south west);
  \draw[flow] (gate.east) -- (111,-151) -- (mix.south);
  \node[anchor=west] at (114,-145) {混合权重};
  \draw[flow] (mix.east) -- (task.west);
\end{tikzpicture}

  \caption{共享载体与耦合阶段的交叉示意。各面板标题以“载体×阶段”标注，上排均使用类别、病害和定位任务，省略重复的图像输入；实线箭头表示前向计算，虚线表示先验联系或训练约束。先验面板的$\bm w_k$仅表示同维、同坐标的参数块，异构输出头省略。输出关系面板画的是训练中的一致性约束；预测传递还会增加真实前向依赖。下排展开条件混合：三个任务各设门控，对同一输入产生各自权重，混合同一组专家输出，再进入任务头。训练时，各任务损失更新本任务的头、门控和所使用的专家；图中省略反向传播路径。标准稠密门控不等于只计算一个专家。}
  \label{fig:ket-sharing-locations}
\end{figure*}
\FloatBarrier

\section{联合学习目标的协调}
\label{sec:ket-sharing-objectives}

现在固定任务组织、共享结构与采样规则，只考察共享参数的一次更新。三个损失相加并不意味着三个任务得到同等照顾：损失数值可能有不同单位，梯度范数可能不同，方向也可能相互抵消。损失大不一定梯度大，梯度大不一定更重要，梯度方向冲突也不等于最终发生负迁移。协调方法必须说明自己处理的是哪一个问题。

一类方法先给任务损失加权，再对标量目标求梯度；另一类方法直接根据各任务梯度构造更新方向。后者也可能求出组合系数，因此两类操作并不互斥。若当前系数相同且求梯度时视其为常量，加权损失与加权梯度可以给出相同更新；区别在于系数由什么依据确定、怎样变化，以及方法试图保证什么性质。

\subsection{任务权重的确定}
\label{sec:ket-sharing-weights}

在同一批次包含所需监督的简化情形下，联合目标写为
\begin{equation}
  L_{\mathrm{joint}}
  =\sum_{k=1}^{K}w_kL_k,
  \qquad w_k\geq0.
  \label{eq:ket-sharing-weighted-loss}
\end{equation}
固定权重时，共享梯度就是$\sum_k w_k\bm g_k$。权重可以表达优先级，也可以校正损失尺度；这两种用途应分开说明。若把定位坐标从米改为厘米，而仍直接计算平方误差，对应损失及其关于同一物理预测参数的梯度会放大$10^4$倍。类别交叉熵没有随之改变，任务的重要性却并未变化。先确定输出归一化、按样本还是按像素平均，再选择任务权重，才能避免把单位变换误当作收益取舍。

共同交付的任务应各有最低要求；主辅学习可以只把主任务作为最终效用，辅助权重通过主任务验证选择；个性化学习还应约定各地区任务是否等权，还是按服务量加权。固定权重的好处是目标清楚、容易复现，困难是同一个权重可能随训练阶段变得失衡。动态权重能适应状态，却不能替用户决定哪项能力值得保护。

检验专用协调器的贡献，需要保留同结构的简单联合训练对照。等权损失提供起点，再在预先限定的搜索预算内扫描固定权重或任务采样比例；各方案使用相同数据、共享结构、基础优化器和停止规则。单任务对照同时改变了共享方式，不能独自隔离协调器的作用。搜索费用与逐任务梯度、额外求解的训练开销都应计入比较；采样比例怎样改变有效目标见第~\ref{sec:ket-sharing-sampling}节。

\term{动态权重平均}（Dynamic Weight Average, DWA）依据各任务近期损失下降速度调权。令
\[
  r_k^{(t-1)}=\frac{L_k^{(t-1)}}{L_k^{(t-2)}},
  \qquad
  w_k^{(t)}
  =K\frac{\exp(r_k^{(t-1)}/T)}
          {\sum_j\exp(r_j^{(t-1)}/T)},
\]
其中$T>0$控制权重差异，前两个阶段须使用预定初始化。下降较慢的任务具有较大的$r_k$，因而获得较大权重；但损失尺度突变、噪声或接近零的分母都会干扰比值。DWA只读取损失轨迹，不表达业务优先级，也不处理梯度方向冲突\scite{ket-liu2019-mtan}

\term{不确定性加权}（uncertainty weighting）从观测模型出发，用任务的噪声尺度解释相对权重。以$d_k$维回归输出为例，设
\begin{equation}
  y_k=f_k(\bm x)+\bm\varepsilon_k,
  \qquad
  \bm\varepsilon_k\sim
  \mathcal N(\bm0,\sigma_k^2\bm I).
\end{equation}
这里$\sigma_k$是任务级同方差噪声尺度，即同一任务对所有输入使用相同的噪声参数，不是每张图像的任意置信分数。定义批次平均平方误差的一半为$L_k^{\mathrm{reg}}$，令$s_k=\log\sigma_k^2$，平均负对数似然为
\begin{equation}
  \mathcal L_k^{\mathrm{reg}}
  =e^{-s_k}L_k^{\mathrm{reg}}
   +\frac{d_k}{2}s_k+C.
  \label{eq:ket-sharing-uncertainty-regression}
\end{equation}
第一项让噪声尺度较大的任务具有较小权重，第二项惩罚无限增大噪声，防止仅靠忽略任务来减小目标。系数$d_k/2$来自输出维数；若又把误差按维度平均，就要一致地调整整项似然，不能只改其中一项。各任务样本数不同且使用平均损失时，这也是明确的任务归一化选择，不再自动等于所有观测联合似然的原始求和。

对分类任务，若温度化的类别概率为$\operatorname{softmax}(e^{-s_k}\bm f_k)$，精确的单样本负对数似然满足
\begin{align}
  \mathcal L_k^{\mathrm{cls}}
  ={}&e^{-s_k}L_k^{\mathrm{CE}}
  +\log\sum_c\exp(e^{-s_k}f_{k,c})
  \nonumber\\
  &-e^{-s_k}\log\sum_c\exp(f_{k,c}),
  \label{eq:ket-sharing-uncertainty-classification}
\end{align}
其中$L_k^{\mathrm{CE}}$是未加温度的交叉熵。Kendall等采用配分函数近似，得到常用的$e^{-s_k}L_k^{\mathrm{CE}}+\tfrac12s_k$形式；它不是任意类别分数下都成立的精确恒等式，不能把回归似然的精确推导直接套给分类\scite{kendall2018}

实际训练可以初始化$s_k=0$，每批先计算各任务损失，再按选定的似然形式同时更新模型参数和$s_k$。对形如$e^{-s_k}L_k+c_ks_k$的一项，有
\begin{equation}
  \frac{\partial\mathcal L_k}{\partial s_k}
  =-e^{-s_k}L_k+c_k.
\end{equation}
固定模型且$L_k>0$时，驻点满足$\sigma_k^2=L_k/c_k$，显示权重与当前残差尺度相联系，而非独立的业务优先级。若训练误差趋近零，噪声估计可能趋向退化边界，需要数值下界或先验等额外约束。噪声假设不合适、困难任务被过早降权，以及稀缺群体被平均噪声掩盖，都是应检查的风险。

GradNorm从另一个问题出发：即使损失已归一化，各任务对共享参数的作用仍可能相差很大。选定一个共享参数块$\bm{\theta}_{\mathrm b}$，通常取靠近任务分支的共享层，定义加权梯度范数
\begin{equation}
  G_k^{(t)}
  =\left\|\nabla_{\bm{\theta}_{\mathrm b}}
     \bigl(w_k^{(t)}L_k^{(t)}\bigr)\right\|_2
  =w_k^{(t)}\|\bm g_{k,\mathrm b}^{(t)}\|_2,
  \label{eq:ket-sharing-gradnorm-size}
\end{equation}
其中权重取正。实现时可令$w_k=\exp(v_k)$或每步投影到正数域，不能依赖一次无约束梯度更新后“恰好仍为正”。再用初始损失归一化学习进度：
\begin{equation}
  q_k^{(t)}=\frac{L_k^{(t)}}{L_k^{(0)}},
  \qquad
  r_k^{(t)}
  =\frac{q_k^{(t)}}{K^{-1}\sum_jq_j^{(t)}}.
  \label{eq:ket-sharing-gradnorm-rate}
\end{equation}
$r_k>1$表示该任务的损失相对自身起点下降得较慢，并不表示它的绝对损失一定更大。若初始损失为零或极小，就需要另设稳定的参照，不能机械相除。

令$\bar G=K^{-1}\sum_kG_k$，GradNorm把任务$k$的目标范数设为$\bar G r_k^\alpha$，其中$\alpha\geq0$控制照顾学习较慢任务的强度，再优化
\begin{equation}
  L_{\mathrm{grad}}
  =\sum_{k=1}^{K}
  \left|G_k-
    \operatorname{stopgrad}(\bar G r_k^\alpha)\right|.
  \label{eq:ket-sharing-gradnorm-objective}
\end{equation}
停止梯度表示在这次权重更新中把目标值视为常量；否则目标本身会跟随权重变化，偏离预定的平衡要求。权重更新时还固定已求得的任务梯度，仅对$w_k$求导；模型参数则用加权任务损失更新，不用$L_{\mathrm{grad}}$直接训练表示。每次权重更新后保持其为正，并作$w_k\leftarrow Kw_k/\sum_jw_j$的归一化。GradNorm以这种方式同时参考梯度大小与相对学习速度\scite{ket-chen2018-gradnorm}

例如，教学设定中三个任务的$q_k$为$(0.2,0.5,0.8)$，则$r_k=(0.4,1,1.6)$。取$\bar G=2,\alpha=1$，目标范数为$(0.8,2,3.2)$：下降较慢的定位任务得到较大的梯度目标。这里只改变权重的更新依据，没有证明定位任务本来就应该更重要。任务因标签噪声而无法继续拟合时，追赶相同训练速度反而可能浪费预算；测量块改变、批次噪声和损失初值也会影响结果。不确定性加权解释噪声尺度，GradNorm解释训练状态，两者都不能免去各任务验证。

\subsection{联合更新方向的构造}
\label{sec:ket-sharing-directions}

即使梯度大小相近，方向仍可能冲突。以下省略阶段上标，固定所有私有参数，将共享更新写成$\bm{\theta}_{\mathrm s}^{+}=\bm{\theta}_{\mathrm s}-\eta\bm d$。对梯度局部Lipschitz连续的任务损失，
\begin{equation}
  L_k(\bm{\theta}_{\mathrm s}-\eta\bm d)
  =L_k(\bm{\theta}_{\mathrm s})
   -\eta\bm g_k^{\top}\bm d+O(\eta^2).
  \label{eq:ket-sharing-direction-effect}
\end{equation}
因此，$\bm g_k^{\top}\bm d>0$意味着足够小的这次更新能降低任务$k$的局部损失。下面分别考察消除冲突分量、寻找共同下降方向，以及在可行方向中分配收益；这些要求可以组合，不是穷尽所有方法的互斥分类。

\subsubsection{基于冲突消解的方向构造}
\label{sec:ket-sharing-pcgrad}

PCGrad从各任务原始共享梯度出发，为每项建立副本$\widetilde{\bm g}_k=\bm g_k$。遍历其他任务$j$时，若当前副本与原始$\bm g_j$内积为负，就执行
\begin{equation}
  \widetilde{\bm g}_k\leftarrow
  \widetilde{\bm g}_k-
  \frac{\widetilde{\bm g}_k^{\top}\bm g_j}
       {\|\bm g_j\|_2^2}\bm g_j.
  \label{eq:ket-sharing-pcgrad-projection}
\end{equation}
被减去的是在$\bm g_j$方向上的冲突分量，投影后与该方向的内积为零。无冲突时保持不变，零范数梯度跳过。其他任务按随机顺序遍历，最后汇总各副本；任务私有参数仍可由本任务损失更新。这里判断冲突用的是已经修正的副本，但投影参照使用其他任务的原始梯度，不能把两个角色混淆\scite{ket-yu2020-pcgrad}

用同一组二维教学损失比较后续三种方法。令$\bm\theta=(\theta_1,\theta_2)^{\top}$为两个共享自由度，取
\begin{align}
  L_1(\bm\theta)
    &=\tfrac12\|\bm\theta-(-1,0)^{\top}\|_2^2,
  \nonumber\\
  L_2(\bm\theta)
    &=\tfrac12\|\bm\theta-(2,-2)^{\top}\|_2^2.
  \label{eq:ket-sharing-toy-losses}
\end{align}
它们只模拟两个任务在共享块上的竞争，不是花卉实验结果。取$\bm\theta=\bm0$，有$\bm g_1=(1,0)^{\top}$、$\bm g_2=(-2,2)^{\top}$，内积为$-2$。原始求和得到$\bm d_{\mathrm{sum}}=(-1,2)^{\top}$，对任务1有$\bm g_1^{\top}\bm d_{\mathrm{sum}}=-1$，因此这次总梯度会使任务1局部变差。

按式~\eqref{eq:ket-sharing-pcgrad-projection}分别投影，
\begin{align}
  \widetilde{\bm g}_1
    &=(1,0)^{\top}-\frac{-2}{8}(-2,2)^{\top}
      =(0.5,0.5)^{\top},
  \nonumber\\
  \widetilde{\bm g}_2
    &=(-2,2)^{\top}-\frac{-2}{1}(1,0)^{\top}
      =(0,2)^{\top}.
\end{align}
于是$\bm d_{\mathrm{PC}}=(0.5,2.5)^{\top}$，与两项原始梯度的内积分别为$0.5$和$4$。取共同步长$\eta=0.1$，实际损失对为
\begin{align}
  (L_1,L_2)(\bm0)&=(0.5,4),
  \nonumber\\
  (L_1,L_2)(-\eta\bm d_{\mathrm{sum}})
    &=(0.625,3.425),
  \nonumber\\
  (L_1,L_2)(-\eta\bm d_{\mathrm{PC}})
    &=(0.4825,3.6325).
  \label{eq:ket-sharing-pcgrad-numbers}
\end{align}
原始求和虽然使总损失下降，却提高了任务1的损失；投影后两者都下降。这只是此处梯度和步长下的结果，两种方向的长度也不同，并不是对所有问题的性能排序。

任务多于两个时，后续投影可能改变前面已经形成的角度，结果会依赖处理顺序；不能把逐对处理说成一次满足全部约束的精确投影。小批次中的负内积也可能只是梯度估计噪声。PCGrad需要分别取得任务梯度并执行额外内积，改变的是局部优化行为，不保证最终每项任务都优于同等预算的单任务模型。

\subsubsection{基于共同下降的方向构造}
\label{sec:ket-sharing-mgda}

如果要求本步对所有当前任务都下降，可以直接寻找这样的方向。\term{多梯度下降算法}（multiple gradient descent algorithm, MGDA）考察任务梯度的最小范数凸组合：
\begin{align}
  \bm\alpha^\star
  &=\argmin_{\substack{\alpha_k\geq0\\\sum_k\alpha_k=1}}
  \frac12\left\|\sum_{k=1}^{K}\alpha_k\bm g_k\right\|_2^2,
  \nonumber\\
  \bm d^\star&=\sum_k\alpha_k^\star\bm g_k.
  \label{eq:ket-sharing-mgda}
\end{align}
它在所有凸组合中找离原点最近的向量，不是事先假定某个任务具有较大业务权重。多任务多目标优化将这一构造用于共享网络，并讨论了降低求解成本的近似\scite{ket-sener2018-mgda}

为何最小范数与共同下降有关？从最优组合朝任何一个顶点$\bm g_k$移动，都不能继续减小平方范数，因此
\begin{equation}
  (\bm g_k-\bm d^\star)^{\top}\bm d^\star\geq0,
  \qquad
  \bm g_k^{\top}\bm d^\star
  \geq\|\bm d^\star\|_2^2.
  \label{eq:ket-sharing-common-descent}
\end{equation}
若$\bm d^\star\ne\bm0$，每个内积均为正，故$-\bm d^\star$是共同的局部下降方向。若各梯度局部具有光滑常数$\beta_k$，下降引理进一步给出
\begin{equation}
  L_k(\bm\theta-\eta\bm d)
  \leq L_k(\bm\theta)
  -\eta\bm g_k^{\top}\bm d
  +\frac{\beta_k\eta^2}{2}\|\bm d\|_2^2.
\end{equation}
步长必须足够小，才能把一阶结论变为实际损失下降；非零方向本身不能保证任意学习率都合适。光滑性、凸性及相应优化基础见第~\ref{chap:convex-learning-and-stability}章。

对式~\eqref{eq:ket-sharing-toy-losses}在零点的梯度，令$\alpha_1=a,\alpha_2=1-a$，凸组合为$(3a-2,2-2a)^{\top}$。最小化其平方范数得到$a=10/13$，因此
\begin{equation}
  \bm d^\star=(4/13,6/13)^{\top},
  \qquad
  \bm g_1^{\top}\bm d^\star
  =\bm g_2^{\top}\bm d^\star
  =4/13.
  \label{eq:ket-sharing-mgda-numbers}
\end{equation}
这次确实找到公共下降方向。若移动到两损失中心的中点$\bm\theta=(0.5,-1)^{\top}$，两梯度变成$(1.5,-1)^{\top}$和$(-1.5,1)^{\top}$，等权凸组合为零。它们要求相反的变化，已经没有对两项损失都严格下降的一阶方向。

一般地，$\bm d^\star=\bm0$表示零向量属于梯度的凸包，这称为此处的\term{帕累托驻点}（Pareto stationary point）：不存在同时严格降低各目标的一阶方向。它不同于\term{帕累托有效}（Pareto efficient），后者要求不存在另一个可行解在所有目标上都不差、且至少一项更好。非凸问题的驻点还可能是较差的局部解；在上述凸二次例子中，中点恰好是有效取舍，但这个例子的性质不能推广到任意网络。若这里只计算共享块梯度，结论也只针对固定私有参数的共享块，任务头仍可能继续改善。两种性质都不说明结果优于各任务自己的单任务基线。

任务较多时，可以在单纯形上用Frank--Wolfe迭代近似求解。给定$\bm\alpha$及$\bm d=\sum_k\alpha_k\bm g_k$，选择使$\bm g_j^{\top}\bm d$最小的任务$j$，再沿$\bm\alpha\leftarrow(1-\gamma)\bm\alpha+\gamma\bm e_j$做一维步长搜索，其中$\bm e_j$是单纯形顶点。一个可检查的停止量为
\begin{equation}
  \Delta=\|\bm d\|_2^2-\min_k\bm g_k^{\top}\bm d.
\end{equation}
若近似解满足$\Delta<\|\bm d\|_2^2$，所有真实任务梯度与$\bm d$的内积仍为正；仅报告求解器运行了几步，却不检查误差，不能继承精确解的共同下降结论。

直接保存全部共享梯度可能昂贵，因此MGDA-UB在共享表示上构造较便宜的上界问题。设$\bm z=h_{\bm{\theta}_{\mathrm s}}(\bm x)$、$\bm r_k=\nabla_{\bm z}L_k$，雅可比为$\bm J=\partial\bm z/\partial\bm{\theta}_{\mathrm s}$，则$\bm g_k=\bm J^{\top}\bm r_k$。最小化$\|\sum_k\alpha_k\bm r_k\|_2$，不自动等于最小化真实参数梯度范数；真实内积为
\begin{equation}
  \bm g_k^{\top}\bm d
  =\bm r_k^{\top}\bm J\bm J^{\top}
    \sum_j\alpha_j\bm r_j.
  \label{eq:ket-sharing-mgda-feature-metric}
\end{equation}
对任意凸组合$\bm r(\bm\alpha)=\sum_k\alpha_k\bm r_k$，谱范数给出
\[
  \left\|\sum_k\alpha_k\bm g_k\right\|_2
  =\|\bm J^\top\bm r(\bm\alpha)\|_2
  \leq\|\bm J\|_2\,\|\bm r(\bm\alpha)\|_2.
\]
这就是以表示梯度范数控制参数梯度范数的上界；$\|\bm J\|_2$对同一参数状态和批次不随$\bm\alpha$改变，但最小化上界不保证最小化真实范数，也不直接给出每项真实内积为正。因此，特征空间内积与参数空间内积可能不同。例如，$\bm J\bm J^{\top}$在这些特征梯度张成的子空间上等于$c\bm I$且$c>0$，是保留该下降关系的一种充分条件；一般情形需核对近似方法的条件，或直接复查真实梯度内积。批次梯度的共同下降也只针对当前批次，不是总体风险逐步下降的保证。

\subsubsection{基于协商收益的方向构造}
\label{sec:ket-sharing-nash}

共同下降可能有许多方向，它们给各任务带来的改善不同。Nash-MTL把选择问题写成任务间的协商：以不更新为协商破裂时的参考状态，把一次小步的一阶下降量作为收益，再在有界方向集合中选择联合收益\scite{ket-navon2022-nash}固定方向半径为1，省略所有任务共同的步长因子，任务效用为$u_k(\bm d)=\bm g_k^{\top}\bm d$，协商目标是
\begin{equation}
  \max_{\substack{\|\bm d\|_2\leq1\\
                  \bm g_k^{\top}\bm d>0\ \forall k}}
  \sum_{k=1}^{K}\log(\bm g_k^{\top}\bm d).
  \label{eq:ket-sharing-nash-objective}
\end{equation}
正效用约束意味着每个任务相对零更新都应得到局部收益。对数和等价于最大化效用乘积，既奖励总的改善，也不容许靠牺牲某项任务为负效用来补偿其他任务。它仍是在给定损失和局部线性近似下协商，不是直接最大化最终验证指标。

令$\bm G=[\bm g_1,\ldots,\bm g_K]$，$\bm H=\bm G^{\top}\bm G$。在存在严格共同下降方向时，最优方向位于球面。对球约束使用拉格朗日乘子$\lambda$，方向的驻点条件为
\begin{equation}
  \sum_k\frac{\bm g_k}{\bm g_k^{\top}\bm d}
  =\lambda\bm d.
\end{equation}
两边与单位向量$\bm d$取内积，得到$\lambda=K$。定义$a_k=1/(\sqrt K\,\bm g_k^{\top}\bm d)$，便有$\bm d=\bm G\bm a/\sqrt K$；代回各任务效用，得到需要求解的正向量方程
\begin{equation}
  \bm H\bm a=\bm a^{-1},
  \qquad
  (\bm a^{-1})_k=1/a_k.
  \label{eq:ket-sharing-nash-system}
\end{equation}
由此$\|\bm G\bm a\|_2^2=K$，单位方向为$\bm d_{\mathrm N}=\bm G\bm a/\sqrt K$，各效用为$1/(\sqrt K a_k)$。这说明系数不仅用于加权，还要共同满足协商的收益条件。非驻点处梯度线性独立使$\bm H$正定，是该求解设定中常用的充分条件；接近相关、秩亏或没有严格共同下降方向时，不能盲目求逆。

原论文用凸代理和凸凹过程近似求解权重。这里另给出便于复算的教学求法，不复现其具体实现。在$\bm H\succ0$时，可最小化下列严格凸函数，其驻点恰好满足式~\eqref{eq:ket-sharing-nash-system}：
\begin{equation}
  \psi(\bm a)
  =\tfrac12\bm a^{\top}\bm H\bm a
   -\sum_k\log a_k,
  \qquad a_k>0.
  \label{eq:ket-sharing-nash-newton}
\end{equation}
其梯度为$\bm r=\bm H\bm a-\bm a^{-1}$，海森矩阵为$\bm H+\operatorname{diag}(a_k^{-2})$。从全正初值出发，解线性方程
\[
  \bigl[\bm H+\operatorname{diag}(a_k^{-2})\bigr]
  \bm\delta=-\bm r
\]
得到牛顿增量。回溯缩短步长，使更新后的所有分量保持为正且$\psi$下降；残差达到预设容差后，再将$\bm G\bm a$归一化为单位方向。

每次重估都要取得任务梯度、形成Gram矩阵并解任务规模的子问题。间隔若干训练步复用权重可减少开销，但权重将不再精确对应当前梯度。

仍用零点处的两个教学梯度，$\bm H=\left[\begin{smallmatrix}1&-2\\-2&8\end{smallmatrix}\right]$。方程为$a_1(a_1-2a_2)=1$和$a_2(-2a_1+8a_2)=1$，两式相减得$a_1^2=8a_2^2$。取正根，
\begin{align}
  a_2&=(8-4\sqrt2)^{-1/2},
  \qquad a_1=2\sqrt2\,a_2,
  \nonumber\\
  \bm d_{\mathrm N}&\approx(0.3827,0.9239)^{\top}.
\end{align}
为了只比较收益分配，将MGDA方向也归一化为$\bm d_{\mathrm M}=(2,3)^{\top}/\sqrt{13}$。两种单位方向的效用分别为
\begin{align}
  (u_1,u_2)(\bm d_{\mathrm M})
    &\approx(0.5547,0.5547),
  \nonumber\\
  (u_1,u_2)(\bm d_{\mathrm N})
    &\approx(0.3827,1.0824).
  \label{eq:ket-sharing-nash-numbers}
\end{align}
MGDA在这个例子中给出相同的一阶收益；Nash方向让任务1取得较小但仍为正的收益，使效用乘积由$4/13\approx0.3077$增至$\sqrt2-1\approx0.4142$。它选择的是另一种取舍，不能据此宣称每项收益都超过MGDA。

图~\ref{fig:ket-sharing-gradient-geometry}把上述计算放在同一组二维坐标中：PCGrad修改冲突分量，MGDA选择梯度凸包中离原点最近的点，Nash则按效用乘积选择单位方向。图中画的是梯度及其组合$\bm d$，实际参数位移为$-\eta\bm d$，不可沿图中箭头直接读取参数移动方向。

\begin{figure*}[htbp]
  \centering
  % Teaching geometry from eq:ket-sharing-toy-losses; no experimental data.
% Source gradients are g1=(1,0), g2=(-2,2), not (-2,3).
% All panels use 11.5 mm per coordinate unit; canvas is 169 x 74 mm.
% Arrows show gradients or their combinations d; displacement is -eta*d.
\begin{tikzpicture}[
  x=1mm,y=1mm,
  font=\figurefont\footnotesize,
  text=BookInk,
  >={Latex[length=1.7mm,width=1.2mm]},
  axis/.style={draw=BookMuted,line width=0.6pt,-{Latex}},
  tick/.style={draw=BookMuted,line width=0.5pt},
  original/.style={draw=FigureInput,line width=0.9pt,-{Latex}},
  projected/.style={draw=FigureModel,line width=1pt,-{Latex}},
  result/.style={draw=FigureLoss,line width=1.1pt,-{Latex}},
  guide/.style={draw=BookMuted,dashed,line width=0.7pt},
  leader/.style={draw=BookMuted,line width=0.5pt},
  heading/.style={anchor=west,font=\figurefont\footnotesize\bfseries},
  label/.style={inner sep=0.6mm},
  note/.style={align=center,inner sep=0pt}
]
  \path[use as bounding box] (0,0) rectangle (169,74);
  \node[heading] at (0,70) {(a) PCGrad：投影};
  \node[heading] at (58,70) {(b) MGDA：凸包最短点};
  \node[heading] at (116,70) {(c) 单位方向与共同下降};

  % Identical origins, ranges and physical scales in all three panels.
  \foreach \offset in {0,58,116} {
    \begin{scope}[shift={(\offset,0)}]
      \begin{scope}[shift={(29,25)},x=11.5mm,y=11.5mm]
        \draw[axis] (-2.25,0) -- (1.45,0);
        \draw[axis] (0,-1.2) -- (0,2.9);
        \foreach \v in {-2,-1,1} {
          \draw[tick] (\v,-0.045) -- (\v,0.045);
          \node[label,anchor=north] at (\v,-0.09) {$\v$};
        }
        \foreach \v in {-1,1,2} {
          \draw[tick] (-0.045,\v) -- (0.045,\v);
          \node[label,anchor=east] at (-0.09,\v) {$\v$};
        }
        \node[label,anchor=north east] at (-0.06,-0.06) {$0$};
      \end{scope}
      \node[label,anchor=west] at (46.5,25) {$d_1$};
      \node[label,anchor=south] at (29,60) {$d_2$};
    \end{scope}
  }

  \begin{scope}[shift={(29,25)},x=11.5mm,y=11.5mm]
    \coordinate (p1) at (0.5,0.5);
    \coordinate (p2) at (0,2);
    \draw[guide] (1,0) -- (p1);
    \draw[guide] (-2,2) -- (p2);
    \draw[original] (0,0) -- (1,0);
    \draw[original] (0,0) -- (-2,2);
    \draw[projected] (0,0) -- (p1);
    \draw[projected] (0,0) -- (p2);
    \draw[result] (0,0) -- (0.5,2.5);
    % The two removed components are orthogonal to the projected vectors.
    \draw[leader] (0.42,0.42) -- (0.50,0.34) -- (0.58,0.42);
    \draw[leader] (-0.14,2) -- (-0.14,1.86) -- (0,1.86);
  \end{scope}
  \node[label,anchor=west] at (42,29) {$\bm g_1$};
  \node[label,anchor=south] at (6,49) {$\bm g_2$};
  \node[label,anchor=west] at (38,35) {$\widetilde{\bm g}_1$};
  \draw[leader] (p1) -- (38,34);
  \node[label,anchor=south] at (21,54) {$\widetilde{\bm g}_2$};
  \draw[leader] (p2) -- (21,54);
  \node[label,anchor=west] at (37,56) {$\bm d_{\mathrm{PC}}$};
  \node[note] at (26.5,4) {$\bm d_{\mathrm{PC}}=(0.5,2.5)^{\top}$};

  \begin{scope}[shift={(87,25)},x=11.5mm,y=11.5mm]
    \coordinate (mgda) at ({4/13},{6/13});
    \draw[draw=BookInk,line width=1pt] (-2,2) -- (1,0);
    \fill[FigureInput] (-2,2) circle[radius=0.9mm];
    \fill[FigureInput] (1,0) circle[radius=0.9mm];
    \draw[projected] (0,0) -- (mgda);
    \fill[FigureModel] (mgda) circle[radius=0.7mm];
    % Perpendicularity to g1-g2=(3,-2) identifies the closest point.
    \draw[leader]
      ({4/13+0.16*3/sqrt(13)},{6/13-0.16*2/sqrt(13)})
      -- ({4/13+0.16/sqrt(13)},{6/13-0.16*5/sqrt(13)})
      -- ({4/13-0.16*2/sqrt(13)},{6/13-0.16*3/sqrt(13)});
  \end{scope}
  \node[label,anchor=south] at (64,49) {$\bm g_2$};
  \node[label,anchor=west] at (99,29) {$\bm g_1$};
  \node[label,anchor=west] at (96,37) {$\bm d^\star$};
  \draw[leader] (mgda) -- (96,36);
  \node[note] at (84.5,4) {$\bm d^\star=(4/13,6/13)^{\top}$};

  \begin{scope}[shift={(145,25)},x=11.5mm,y=11.5mm]
    % g1.d>0 and g2.d>0 give d1>0 and d2>d1: angles 45 to 90 deg.
    \fill[FigureModel!12] (0,0) -- (45:1)
      arc[start angle=45,end angle=90,radius=11.5mm] -- cycle;
    \draw[draw=BookMuted,dotted,line width=0.7pt] (0,0) circle[radius=11.5mm];
    \draw[guide] (45:1) -- (0,0) -- (90:1);
    \coordinate (unitm) at ({2/sqrt(13)},{3/sqrt(13)});
    \coordinate (unitn) at ({sin(22.5)},{cos(22.5)});
    \draw[projected] (0,0) -- (unitm);
    \draw[result,dashed] (0,0) -- (unitn);
  \end{scope}
  \node[note,anchor=north west,align=left] at (118,54)
    {共同下降区域\\$d_1>0$\\$d_2>d_1$};
  \node[label,anchor=south] at (151,44) {$\bm d_{\mathrm N}$};
  \draw[leader] (unitn) -- (151,44);
  \node[label,anchor=west] at (159,38) {$\bm d_{\mathrm M}$};
  \draw[leader] (unitm) -- (159,38);
  \node[label,anchor=east] at (133,12) {单位圆};
  \draw[leader] (134,13) -- (138.1,15.8);
  \node[note] at (142.5,4) {实际位移：$-\eta\bm d$};
\end{tikzpicture}

  \caption{同一二维教学算例中的梯度协调。三面板使用相同坐标尺度，原始梯度为$\bm g_1=(1,0)^{\top}$、$\bm g_2=(-2,2)^{\top}$。(a)虚线连接原始端点与PCGrad投影端点，投影后的两向量相加得到$\bm d_{\mathrm{PC}}$。(b)线段是原始梯度的凸包，最近点为$\bm d^\star=(4/13,6/13)^{\top}$。(c)点线圆为单位圆，浅色扇区表示单位圆内满足$d_1>0,d_2>d_1$的组合，即$-\bm d$使两任务一阶共同下降；虚线径向边界不含在严格下降区域内。实线$\bm d_{\mathrm M}$为归一化MGDA方向，虚线$\bm d_{\mathrm N}$为Nash方向。图中均画梯度组合，实际位移是$-\eta\bm d$；局部几何不构成最终性能排名。}
  \label{fig:ket-sharing-gradient-geometry}
\end{figure*}
\FloatBarrier

在同一可行集和零参考效用下，最大化对数和具有比例公平的含义：相对于最优效用，其他可行方案的相对收益变化之和不为正。把某个任务损失乘以正数，只给其对数效用增加常量，也不会改变这里的最优单位方向。但这种尺度性质不替代任务优先级，更不等于不同人群得到公平服务。存在群体代价时，仍需选择相应评价单位与约束；任务总损失下降、平均指标提高和每项任务均受益应分别报告。

语言任务的实证对照进一步说明了为何要调好简单联合训练基线。Xin等在英语译为法语、并分别搭配汉语、德语或罗马尼亚语的联合翻译实验中，比较了PCGrad、MGDA、GradNorm、IMTL-G及随机加权方法。所测专用方法没有超出简单任务采样比例扫描得到的性能取舍，方法之间的排序还随评价时的任务权重变化\scite{ket-xin2022-mto}这一负面结果限于该研究的任务、模型和调参范围；其中未比较Nash-MTL，不能据此否定它或所有梯度协调方法。

这些协调方法针对训练中已见的任务。若要证明共享结构能够帮助从未共同训练过的新任务，应采用第~\ref{sec:ket-transfer-task-splits}节的任务层留出与适应设定，而不是只留出已见任务的若干样本。持续冲突还可能表明任务分组、共享深度或输出约束不合适；优化器只能在当前结构内安排更新，不能凭空恢复已经被结构排除的任务信息。

\section{联合训练资源的分配}
\label{sec:ket-sharing-resources}

联合目标确定以后，还需使训练过程真正实现它。数据与标签决定能获得什么证据，采样决定这些证据被使用多少次，容量与计算预算决定每次能执行什么模型。只报告总训练步数，无法分辨增加独立标注、重复稀缺样本和扩大网络所带来的不同投入。

\subsection{样本与标签资源的分配}
\label{sec:ket-sharing-labels}

三个任务可以使用完全相同的图像，也可以只有部分对象重叠。在固定标注预算下，给更多独立植株标类别，增加了类别任务的覆盖；给已标类别的图像补充病害和位置，则使多个任务有机会围绕同一对象学习对应关系。均匀分配标签容易管理，按任务缺口分配可能缓解最薄弱的任务，成套标注则便于估计输出关系。哪一种更合适取决于瓶颈，不能把相同标注数直接视为相同信息量。

令$m_{i,k}\in\{0,1\}$表示样本$i$是否具有任务$k$的可用标签。对批次$B$，只在$n_{B,k}=\sum_{i\in B}m_{i,k}>0$时计算
\begin{equation}
  L_k(B)=
  \frac{1}{n_{B,k}}
  \sum_{i\in B}m_{i,k}
  \ell_k(f_k(\bm x_i),y_{i,k}).
  \label{eq:ket-sharing-label-mask}
\end{equation}
实现时只读取掩码为1的标签；数学中的乘零不能成为访问未定义标签的理由。若$n_{B,k}=0$，式~\eqref{eq:ket-sharing-label-mask}没有定义，本批也没有可用的该任务监督项。缺少病害标签的图像不等于健康图像，更不能把未知类别写成负类。

例如，教学批次的第一张图像只有类别标签，第二张同时具有类别、病害和定位标签，第三张只有病害与定位标签。类别的两个可用损失为$0.2,0.6$，病害为$0.8,0.4$，定位为$0.3,0.5$，于是三个任务平均损失为$0.4,0.6,0.4$。若把每项都除以批次大小3，就额外乘上了本批标签可见率$2/3$；在标签数量不等时，这种归一化还会不均匀地改变任务作用。两种归一化各自可以定义目标，但必须明确，而不是由实现细节偶然决定。

标签掩码只删除无监督项，并不纠正非随机缺失。若专家优先标注明显病害，式~\eqref{eq:ket-sharing-label-mask}估计的是被标注风险$\mathbb E[\ell_k\mid m_k=1]$，一般不等于式~\eqref{eq:ket-sharing-task-risk}中的总体风险。要恢复后者，先固定任务$k$和当前预测器，设批次$B$从该任务的目标输入总体独立同分布抽取，再按标注机制观察其中部分标签。

记$\ell_{i,k}=\ell_k(f_k(\bm x_i),y_{i,k})$。假设缺失在给定输入后可以忽略，即$m_{i,k}\perp y_{i,k}\mid\bm x_i$：同样的输入条件下，是否标注不再依赖未知答案。还要求可见概率$e_k(\bm x)=\mathbb P(m_{i,k}=1\mid\bm x_i=\bm x)$正确，且在目标总体的支撑上满足$e_k(\bm x)>0$。在损失可积时，完整批次分母给出
\begin{equation}
  \widehat R_k(B)=\frac{1}{|B|}
    \sum_{i\in B}\frac{m_{i,k}}{e_k(\bm x_i)}\ell_{i,k}.
  \label{eq:ket-sharing-label-ipw}
\end{equation}
它为何估计总体风险，可以逐项检查条件期望：
\begin{align}
  \mathbb E\!\left[
    \frac{m_{i,k}\ell_{i,k}}{e_k(\bm x_i)}
    \,\middle|\,\bm x_i,y_{i,k}\right]
  &=\frac{\ell_{i,k}}{e_k(\bm x_i)}
    \mathbb P(m_{i,k}=1\mid\bm x_i,y_{i,k})
  \nonumber\\
  &=\ell_{i,k}.
  \label{eq:ket-sharing-label-ipw-expectation}
\end{align}
再对输入和标签取期望，即得$\mathbb E[\widehat R_k(B)]=R_k(f_k)$。这里对未见标签的损失只作数学定义，实现仍仅累加可见项。若用估计的$e_k$代入，就要另行考虑概率估计误差，不能沿用严格无偏的结论；若可见性仍依赖给定输入后未知的答案，或某些目标区域$e_k=0$，上述恢复依据便不成立。

完整批次分母不能换成可见标签数。例如，四个样本的损失都为1，可见概率都为$1/2$，本批恰有两个标签，则加权和为4，除以$|B|=4$得到1；仍除以$n_{B,k}=2$却得到2。另一种选择是除以权重和，形成\term{自归一化估计}（self-normalized estimator）：
\begin{equation}
  \widehat R_{k,\mathrm{SN}}(B)
  =\frac{\sum_{i\in B}m_{i,k}\ell_{i,k}/e_k(\bm x_i)}
    {\sum_{i\in B}m_{i,k}/e_k(\bm x_i)}.
  \label{eq:ket-sharing-label-self-normalized}
\end{equation}
该式只在分母为正时定义。它使用随机权重和归一化，有限样本下一般有偏；在上面的常损失例子中恰好也得到1，不代表两种估计量总是相同。

无标签批次尤其需要区分三个对象。式~\eqref{eq:ket-sharing-label-ipw}此时取值为0，这是随机估计量的一个取值，不是观测到任务的真实损失为零；可见标签平均与自归一化估计则都未定义。可以省略该任务本批的监督反传，但若仍声称估计原总体，就应把这个零贡献保留在完整批次计数与风险平均中，也不按本批有标签的任务重新归一化预设权重。事实上，在固定模型下只保留$n_{B,k}>0$的批次再平均，会得到条件期望$R_k(f_k)/\mathbb P(n_{B,k}>0)$，一般已不是原风险。

GradNorm的训练速度记录是另一个状态，不能把上述随机零值写成“已经学到零损失”。缺标时保留该任务最近的有效进度记录；一种保守做法是暂缓本步整组GradNorm权重更新，保持上一组权重，监督梯度仍按已经选定的估计式处理。若采用异步进度或只对有标签任务调权，须另行规定更新与归一化规则。重要性校正说明的是固定预测器的风险估计；在可交换求导与期望时才进一步得到无偏梯度，并不自动保证动态调权后的完整更新仍无偏。相关分布校正机制见第~\ref{sec:ket-update-weighting}节。

核算样本还应区分独立对象、原始图像、增强副本与任务标签记录。一张图像有三种标签，提供三类监督，却不是三张独立照片；同一植株的连续拍摄也不能随意分到训练和测试两边。新增标注怎样查询与校验属于知识获取问题，本节负责安排已获标签及其预算，不把多次使用同一标签算成新增证据。

指令任务使这种区别更容易被忽略。T0将不同数据集的任务改写为带自然语言提示的输入输出要求，同一任务可以具有多个模板，并通过留出任务类别检验未参与共同训练的预测要求\scite{ket-sanh2022-t0}因此，应分别记录任务语义、数据集和模板：见过同一问题的另一种问法，考验的是措辞变化；在同类预测要求下更换数据来源，考验数据集变化；留出整个预测类别，才对应另一种任务泛化问题。这三种“未见”不能用一个测试分数代替。

FLAN的扩展研究进一步整理多种指令数据来源，区分数据集、任务类别、任务及模板，并在构造混合时处理来源重叠\scite{ket-chung2024-flan}对花卉系统也应采取同样的核算意识：把“判断是否患病”改写成多条提示，不应让病害任务在数据清单中被重复计算成多个独立能力。合并集合时要检查原始对象、改写文本和答案是否重复，尤其防止训练内容与评价内容重叠；最终报告应说明是增加了预测任务、独立样本，还是表达形式。

\subsection{训练更新机会的分配}
\label{sec:ket-sharing-sampling}

本节先讨论采样前权重已经确定的线性加权损失更新。PCGrad、MGDA与Nash-MTL等多梯度协调方法，需要在同一参数状态下取得参与协调的各任务梯度；若分步收集或使用旧梯度，则须另行规定收集与近似规则。下述逆概率校正可以恢复线性加权梯度的期望，不能直接提供这些非线性协调算子的性质。

假设每步只抽取一个任务，再从该任务的已标注训练集中均匀抽一个样本。令任务抽取概率为$p_k>0$，选中后乘以权重$w_k$。在当前参数、采样概率及权重固定的条件下，随机梯度的期望为
\begin{equation}
  \mathbb E[\widehat{\bm g}]
  =\sum_{k=1}^{K}p_kw_k
    \nabla\hat R_k.
  \label{eq:ket-sharing-effective-objective}
\end{equation}
因此，未作校正时的有效目标是$\sum_kp_kw_k\hat R_k$，而不是仅看代码中写下的$\sum_kw_k\hat R_k$。这个推导假定任务内均匀采样；若任务内也按难度选择样本，还要继续核对条件采样分布。

按样本比例取$p_k=n_k/\sum_jn_j$时，大数据任务得到更多更新，等权损失对应标签记录的样本平均；按任务均衡取$p_k=1/K$时，各任务得到相同的期望机会，等权损失对应任务平均。两者可以都合理，但表达的目标不同。总步数$T$固定时，任务$k$的期望抽取次数为$Tp_k$，并非每次随机训练都恰好等于这个数。提高频率让任务在更多批次上产生梯度；提高每次权重只是放大已抽到的梯度，不会自动增加独立观察。

若预先指定目标为$\sum_kq_k\hat R_k$，其中$q_k\geq0,\sum_kq_k=1$，可以在抽到任务$k$时使用
\begin{equation}
  \widehat{\bm g}
  =\frac{q_k}{p_k}\nabla\ell_k,
  \qquad
  \mathbb E[\widehat{\bm g}]
  =\sum_kq_k\nabla\hat R_k.
  \label{eq:ket-sharing-importance-sampling}
\end{equation}
这里要求$q_k>0$的任务也有$p_k>0$，并把当前$q_k/p_k$作为梯度估计系数，而不是对采样策略本身求导。自适应$p_k$在给定训练历史后满足这些条件时，也可以逐步保持相同的条件期望目标。若原来已经设置损失权重，需检查最终乘积，不能重复校正。

无偏并不等于稳定。记$\bm g_k$为任务内随机样本梯度，$\bm\mu_k=\mathbb E[\bm g_k]$，校正估计的均方波动为
\begin{align}
  &\mathbb E\left\|
    \widehat{\bm g}-\sum_kq_k\bm\mu_k
  \right\|_2^2
  \nonumber\\
  &\quad=
  \sum_k\frac{q_k^2}{p_k}\mathbb E\|\bm g_k\|_2^2
  -\left\|\sum_kq_k\bm\mu_k\right\|_2^2.
  \label{eq:ket-sharing-sampling-variance}
\end{align}
某项$p_k$过小时，偶尔抽到它会产生很大的校正梯度。适当提高其采样概率可以降低这类波动，但若样本池很小，也只是重复使用相同证据；是否过拟合仍需验证。设置概率下界或截断权重可改善数值稳定性，截断则通常引入目标偏差，应明确记录。

训练早期偏向稳定的类别监督，后期增加定位和病害机会，是一种课程安排；依据验证集上哪些任务尚未达到要求来调整频率，是另一种反馈分配。比较时应固定总更新或总计算预算，记录各阶段的实际任务次数以及是否做了重要性校正。验证反馈可以用来决定训练机会，最终测试结果不能被反复用于调整采样。动态采样若未校正，通常也在动态改变目标，不能只称为更高效地求解同一个目标。

\subsection{计算与存储预算的分配}
\label{sec:ket-sharing-budget}

参数共享可能节省部署空间，但节省多少取决于调用方式。设一个骨干有$P_{\mathrm b}$个参数，每个任务头有$P_{\mathrm h}$个参数。三个独立模型共需$3(P_{\mathrm b}+P_{\mathrm h})$个参数；一个共享骨干加三个头需$P_{\mathrm b}+3P_{\mathrm h}$个参数。若同一图像同时请求三项输出，共享骨干还可能只计算一次；若三个请求来自不同图像，就不能复用同一次中间激活。只请求类别时，也不能假定病害与定位的训练费用已经自然摊销。

例如，以下参数量仅为教学构造，以百万参数为单位。取$P_{\mathrm b}=12$、$P_{\mathrm h}=1$，独立模型总计39，共享结构总计15；共享模型单任务路径涉及13。这个比较说明结构如何改变费用，不能证明15的共享模型比39的独立模型更好。若控制相同总参数预算39，共享方案可以把骨干扩为36、三个头各保留1，或把预算用于私有分支，再实际检验任务表现。总参数相同也不必然意味着浮点运算量、时延或显存峰值相同。

条件专家需要额外区分\term{总容量}与\term{激活容量}：前者是保存的全部可用参数，后者是一次请求实际参与计算的参数。设公共前端有$P_0$个参数，每个专家有$P_{\mathrm e}$个参数，任务门控及任务头分别为$P_{\mathrm g,k},P_{\mathrm h,k}$，则
\begin{align}
  P_{\mathrm{total}}
    &=P_0+EP_{\mathrm e}
      +\sum_k(P_{\mathrm g,k}+P_{\mathrm h,k}),
  \nonumber\\
  P_{\mathrm{active},k}
    &=P_0+|A_k(\bm x)|P_{\mathrm e}
      +P_{\mathrm g,k}+P_{\mathrm h,k},
  \label{eq:ket-sharing-expert-budget}
\end{align}
其中$A_k(\bm x)$是实际执行的专家集合。稠密MMoE通常有$|A_k|=E$；稀疏选择才可能让它更小。多个任务共同处理同一输入时，能缓存复用的专家计算由集合$\bigcup_kA_k(\bm x)$决定，而不是简单相加或默认只运行一次。即使不执行某个专家，它的参数也可能仍驻留在设备中。

取$P_0=4,E=4,P_{\mathrm e}=6$，三个任务各有$P_{\mathrm g,k}=0.1,P_{\mathrm h,k}=1$，总参数为31.3。稠密执行一个任务的激活参数为29.1；若另行采用每次只执行一个专家的稀疏机制，则为11.1。总容量仍是31.3，不能把11.1写成模型存储量。门控、跨设备传输、专家输入聚合与批次形状也会影响实际运行成本，因此激活参数只是一项计数，不等于精确计算量或时延。

训练状态还包括梯度和优化器状态。以全部可训练参数使用FP32和Adam、分别保存参数、梯度及两个动量为例，不计激活及其他缓冲时约需$16P$字节。15百万参数约对应240\,MB，而仅保存参数是60\,MB，这里MB按$10^6$字节计。条件专家即使每步只激活部分参数，也通常要为已经训练过的全部参数保留优化器状态；混合精度主副本、重计算、状态分片与卸载会改变核算，不能沿用上述简化数字。训练峰值还受批次大小、图像分辨率和保存的中间激活影响，具体执行与通信优化交给第~\ref{part:systems}部分。

表~\ref{tab:ket-sharing-resource-controls}把三种常被混淆的改变放在同一组任务上。这里假设类别、病害和定位分别具有$600,300,100$条可抽取的标签记录，每步抽一个任务及一条记录，共训练1000步；每行固定数据、优化器与损失单位。$C_0$表示基准共享容量，不预设它与计算量成线性关系。只改变权重与只改变采样，都能形成任务均衡的期望目标，却具有不同的更新频率与方差；增加共享容量则改变假设空间与每次计算，不能归入同一种“提高任务权重”。

\begin{table*}[htbp]
  \centering
  \small
  \begin{tabularx}{\linewidth}{p{18mm}p{36mm}p{25mm}p{28mm}X}
    \toprule
    教学配置 & 概率$\bm p$与权重$\bm w$ & 有效系数$p_kw_k$ & 期望抽取次数 & 容量、代价与控制 \\
    \midrule
    基准 &
    $\bm p=(0.6,0.3,0.1)$\newline
    $\bm w=(1,1,1)$ &
    $(0.6,0.3,0.1)$ &
    $(600,300,100)$ &
    容量$C_0$；样本比例目标，记录基准计算量。 \\
    只改权重 &
    $\bm p$同基准\newline
    $\bm w=(5/9,10/9,10/3)$ &
    $(1/3,1/3,1/3)$ &
    $(600,300,100)$ &
    容量$C_0$；机会不变，稀少任务单次梯度更大。 \\
    只改采样 &
    $\bm p=(1/3,1/3,1/3)$\newline
    $\bm w=(1,1,1)$ &
    $(1/3,1/3,1/3)$ &
    各任务$1000/3$ &
    容量$C_0$；小任务重复更多，期望计算取决于任务单步费用。 \\
    只改容量 &
    $\bm p,\bm w$同基准 &
    $(0.6,0.3,0.1)$ &
    $(600,300,100)$ &
    容量$2C_0$；机会与系数不变，需另计每步成本。 \\
    联合改变 &
    $\bm p=(0.4,0.3,0.3)$\newline
    $\bm w=(5/6,10/9,10/9)$ &
    $(1/3,1/3,1/3)$ &
    $(400,300,300)$ &
    容量$2C_0$；按概率校正目标，分别控制频率和容量。 \\
    \bottomrule
  \end{tabularx}
  \caption{权重、采样与容量的不同作用。所有数字都是可复算的教学配置，不含实测性能；向量依次对应类别、病害和定位。次数为随机采样的期望而非一次运行的保证。共享块每步都可更新，被选任务的专属头才在该步获得直接监督；1000步不等于相同总计算，特别是在任务费用或容量改变时。}
  \label{tab:ket-sharing-resource-controls}
\end{table*}

资源对照最终要与收益对照一起报告。三个任务指标含义不同时，不能直接把类别准确率、病害召回率和定位误差取平均；可以分别给出原始指标，再在明确的基线归一化或效用定义下报告任务平均、兼容指标的样本平均和最弱任务表现。还应列出每项最低要求是否满足，避免平均改善掩盖病害任务退化。完整投入包括标签、独立对象、采样重复、总训练计算、保存参数、训练状态和部署调用条件。第~\ref{chap:ket-evaluation}章进一步统一这些评价口径。

\begin{samepage}
\section{本章小结}
\label{sec:ket-sharing-summary}

类别、病害与定位能否有效共享，取决于四项相互制约的设计。任务组织规定谁与谁共同学习以及哪些收益必须保留；共享位置决定经验通过共同计算、参数先验还是输出关系流动；联合目标协调尺度、方向与收益取舍；资源分配则决定这些要求在有限标签、更新次数与模型容量下怎样实现。共同参数和统一接口提供了共享的可能性，实际收益仍需与预算相当的单任务或分组方案比较。

\end{samepage}

只部署类别任务时，属性辅助目标应以类别收益和额外成本决定去留；三个任务共同交付时，每项退化都应单独记录；多个地区具有个性化要求时，共同中心或专家之外还需保留必要的局部自由。局部梯度冲突可以提示一次更新的问题，但不等于长期负迁移；参数负相关可以提供统计帮助，也不等于损害。面对共享失败，应依次核查任务关系证据是否适用、共享结构是否丢失必要信息、损失尺度与取舍是否合理，以及采样和容量是否实际实现了指定目标，而不是只继续调整一个权重。

这些选择可以相互反馈。调整分组后要重新估计亲和性，增加专家后要重新核算总容量与使用分布，改变采样后要检查有效目标是否变化。每次调整都应保持评价单位、数据条件与成本口径可比较。对已见任务的共同学习收益由本章的设定检验；对未见任务的迁移收益需要任务层留出，关于长期维持和整体证据强度的问题则分别交给知识维护与统一评价。
